\section{주요 사업 및 금융 기능}

한국산업은행의 사업은 일반 상업은행과 달리, 단순한 예금·대출 서비스가 아니라 \textbf{대규모 기업금융과 투자은행(IB) 기능}을 중심으로 구성된다. :contentReference[oaicite:0]{index=0}  

이는 산업은행이 정책금융기관으로서 민간 금융기관이 수행하기 어려운 영역을 담당한다는 점에서 중요한 의미를 가진다.

\subsection{기업금융 (Corporate Finance)}

산업은행의 핵심 사업은 대기업 및 중견기업을 대상으로 하는 기업금융이다.

일반 은행이 단기 대출 중심의 금융을 제공하는 것과 달리, 산업은행은 장기적 관점에서 기업의 성장과 구조개선을 지원하는 금융을 제공한다.

특히 다음과 같은 영역에서 중요한 역할을 수행한다:

\begin{itemize}
    \item 대규모 설비투자 자금 지원
    \item 전략 산업 및 기간산업 금융
    \item 중장기 성장 자금 공급
\end{itemize}

이는 산업은행이 단순한 자금 공급자가 아니라 \textbf{산업 발전의 방향을 형성하는 금융 주체}임을 의미한다.

\subsection{투자은행(IB) 기능}

산업은행은 국내 금융기관 중에서도 강력한 투자은행 기능을 수행하는 기관 중 하나이다. :contentReference[oaicite:1]{index=1}  

주요 IB 업무는 다음과 같다:

\begin{itemize}
    \item 인수합병(M\&A) 자문 및 금융 지원
    \item 기업 구조조정 관련 금융
    \item 자본시장 관련 투자 및 자문
\end{itemize}

특히 대규모 기업 인수나 구조조정 과정에서 산업은행은 자금 제공자이자 조정자로서 핵심적인 역할을 수행한다.

\subsection{프로젝트 파이낸싱 (Project Financing)}

산업은행은 사회기반시설 및 대형 개발사업에 대한 프로젝트 파이낸싱(PF)을 수행한다.

이러한 사업은 초기 투자 규모가 크고 위험이 높기 때문에 민간 금융기관이 단독으로 참여하기 어렵다.  
산업은행은 이러한 프로젝트에 자금을 공급하고 금융 구조를 설계함으로써 사업의 실행 가능성을 높인다.

대표적인 PF 대상은 다음과 같다:

\begin{itemize}
    \item 인프라 건설 (도로, 철도, 발전소 등)
    \item 에너지 및 자원 개발
    \item 대형 산업 프로젝트
\end{itemize}

이는 산업은행이 \textbf{경제의 기반을 형성하는 투자}에 직접적으로 관여하고 있음을 보여준다.

\subsection{구조조정 및 위기기업 금융}

산업은행은 기업 구조조정 과정에서 중심적인 역할을 수행한다.

부실기업이 발생할 경우, 산업은행은 채권단을 구성하거나 주도하여 다음과 같은 기능을 수행한다:

\begin{itemize}
    \item 채무 재조정 및 자본 재구성
    \item 자회사 매각 및 사업 재편
    \item 신규 자금 지원을 통한 회생 지원
\end{itemize}

이 과정에서 산업은행은 단순한 채권자가 아니라 \textbf{기업 구조조정의 조정자(coordinator)} 역할을 수행한다.

\subsection{금융시장 보완 기능}

산업은행은 시장에서 자금 공급이 부족한 영역을 보완하는 역할도 수행한다.

예를 들어:
\begin{itemize}
    \item 경기 침체 시 기업 유동성 지원
    \item 금융시장 불안 시 자금 공급 확대
\end{itemize}

이러한 기능은 산업은행이 단순한 금융기관이 아니라 \textbf{거시경제 안정화 장치}로 작동하고 있음을 보여준다.

\subsection{민간 금융과의 관계}

산업은행의 사업은 민간 금융기관과 보완적 관계를 형성하는 것을 원칙으로 한다.

그러나 현실에서는 다음과 같은 긴장이 존재한다:

\begin{itemize}
    \item 동일 시장에서의 경쟁
    \item 정부 지원에 따른 경쟁 불균형
    \item 민간 금융 위축 가능성
\end{itemize}

따라서 산업은행의 사업 영역 설정은 정책적으로 중요한 문제로 남아 있다.

\subsection{소결}

한국산업은행의 사업은 단순한 금융상품 제공이 아니라, 기업금융, 투자은행 기능, 프로젝트 파이낸싱, 구조조정 등 \textbf{고위험·대규모 금융 영역}을 중심으로 구성된다.

이는 산업은행이 민간 금융이 충분히 담당하지 못하는 영역을 보완하는 동시에, 국가 경제의 구조적 발전을 지원하는 핵심 금융기관임을 의미한다.